\documentclass[11pt]{article}

\usepackage[a4paper,margin=1in]{geometry}
\usepackage{amsmath,amssymb,bm}
\usepackage{physics}
\usepackage{hyperref}
\usepackage{xcolor}
\usepackage{enumitem}
\usepackage{microtype}

\hypersetup{colorlinks=true, linkcolor=blue, citecolor=blue, urlcolor=blue}
\setlength{\parindent}{0pt}
\setlength{\parskip}{0.65em}
\setlist{nosep}

\title{A Compact Bath-Temperature Extension of the LTE Hammer--Rosen Marshak-Wave Model}
\author{Menahem Krief}
\date{\today}

\newcommand{\Tb}{T_b}
\newcommand{\Ts}{T_s}
\newcommand{\xf}{x_F}
\newcommand{\eps}{\epsilon}
\newcommand{\dif}{\,d}
\newcommand{\Jfour}{J_4}
\newcommand{\Jn}{J_n}

\begin{document}
\maketitle

\begin{abstract}
The original LTE Hammer--Rosen (HR) theory assumes that the surface temperature
$\Ts(t)$ is given.  Here we summarize a compact bath-temperature version in
which the input is the incident bath temperature $\Tb(t)$ and the surface
temperature is estimated from a Marshak boundary condition.  The model keeps
the original HR front and profile structure, avoids differentiating the noisy
bath history, and requires only fluence integrals plus one scalar algebraic
solve at each time.
\end{abstract}

\section{LTE diffusion equation and material model}

The LTE material model is
\begin{equation}
 e=fT^\beta\rho^{-\mu},
 \qquad
 \frac{1}{\kappa_R}=gT^\alpha\rho^{-\lambda},
 \label{eq:material}
\end{equation}
where $e$ is the specific material energy, $\kappa_R$ is the Rosseland opacity,
$\rho$ is constant, and $f,\beta,\mu,g,\alpha,\lambda$ are material constants.
The LTE diffusion equation can be written as
\begin{equation}
 \boxed{
 \frac{\partial T^\beta}{\partial t}
 =C\frac{\partial^2 T^n}{\partial x^2}}
 \label{eq:lte_diffusion}
\end{equation}
with
\begin{equation}
 \boxed{n=4+\alpha,\qquad \eps=\frac{\beta}{n}}
 \label{eq:n_eps}
\end{equation}
and
\begin{equation}
 \boxed{
 C=\frac{16g\sigma}{3nf}\rho^{-2+\mu-\lambda}.}
 \label{eq:C}
\end{equation}
Here $\sigma$ is the Stefan--Boltzmann constant in the unit system used for the
HR coefficients.

\section{Original surface-temperature HR formulas}

For a prescribed surface temperature $\Ts(t)$, define
\begin{equation}
 H(t)=\Ts^n(t).
\end{equation}
The first-order HR front formula is
\begin{equation}
 \boxed{
 \xf^2(t)=K_\eps C\,H^{-\eps}(t)\int_0^t H(t')\dif t'}
 \label{eq:surface_hr_front_H}
\end{equation}
where
\begin{equation}
 \boxed{K_\eps=\frac{2+\eps}{1-\eps}.}
 \label{eq:Keps}
\end{equation}
Equivalently, since $H^{-\eps}=\Ts^{-\beta}$,
\begin{equation}
 \boxed{
 \xf^2(t)=K_\eps C\,\Ts^{-\beta}(t)
 \int_0^t \Ts^n(t')\dif t'.}
 \label{eq:surface_hr_front_T}
\end{equation}
The HR material energy per unit area is
\begin{equation}
 \boxed{
 E(t)=f\rho^{1-\mu}(1-\eps)\xf(t)H^\eps(t).}
 \label{eq:HR_energy}
\end{equation}
Combining \eqref{eq:surface_hr_front_H} and \eqref{eq:HR_energy} gives
\begin{equation}
 \boxed{
 E^2(t)=Z_1H^\eps(t)\int_0^tH(t')\dif t'}
 \label{eq:E_squared}
\end{equation}
with
\begin{equation}
 \boxed{
 Z_1=f^2\rho^{2(1-\mu)}(2+\eps)(1-\eps)C.}
 \label{eq:Z1}
\end{equation}

\section{Bath boundary condition}

The bath-driven problem prescribes $\Tb(t)$ rather than $\Ts(t)$.  We use the
Marshak boundary condition in the convention
\begin{equation}
 \boxed{
 F(t)=2\sigma\left[\Tb^4(t)-\Ts^4(t)\right].}
 \label{eq:marshak}
\end{equation}
Energy conservation gives
\begin{equation}
 E(t)=\int_0^tF(t')\dif t'.
 \label{eq:energy_from_flux}
\end{equation}
The exact Cohen--Heizler boundary-temperature method solves
\eqref{eq:E_squared}--\eqref{eq:energy_from_flux} by time marching.  The compact
model below replaces that marching calculation by a fluence approximation.

\section{Fluence-ratio approximation}

Define the instantaneous surface-to-bath ratio
\begin{equation}
 \boxed{
 u(t)=\left[\frac{\Ts(t)}{\Tb(t)}\right]^4,}
 \label{eq:u_def}
\end{equation}
so that
\begin{equation}
 \boxed{
 \Ts(t)=\Tb(t)u^{1/4}(t).}
 \label{eq:Ts_from_u}
\end{equation}
The approximation is that this ratio is slowly varying over the accumulated
fluence history:
\begin{equation}
 \boxed{
 \frac{\Ts(t')}{\Tb(t')}\approx \frac{\Ts(t)}{\Tb(t)},
 \qquad 0<t'<t.}
 \label{eq:fluence_approx}
\end{equation}
Therefore
\begin{equation}
 \int_0^t\Ts^4(t')\dif t'\approx u(t)\int_0^t\Tb^4(t')\dif t',
\end{equation}
and
\begin{equation}
 \int_0^t\Ts^n(t')\dif t'\approx u^{n/4}(t)\int_0^t\Tb^n(t')\dif t'.
\end{equation}

Define the two bath fluences
\begin{equation}
 \boxed{
 \Jfour(t)=\int_0^t\Tb^4(t')\dif t',
 \qquad
 \Jn(t)=\int_0^t\Tb^n(t')\dif t'.}
 \label{eq:fluences}
\end{equation}
Using \eqref{eq:marshak}, \eqref{eq:energy_from_flux}, and the fluence
approximation gives
\begin{equation}
 E(t)\approx 2\sigma\Jfour(t)\left[1-u(t)\right].
 \label{eq:E_fluence}
\end{equation}
Using \eqref{eq:E_squared} gives
\begin{equation}
 E^2(t)\approx Z_1\,u^{(n+\beta)/4}(t)\,\Tb^\beta(t)\,\Jn(t),
 \label{eq:E_HR_fluence}
\end{equation}
where $n\eps=\beta$ has been used.
Equating \eqref{eq:E_fluence} and \eqref{eq:E_HR_fluence} yields the scalar
algebraic equation for $u(t)$:
\begin{equation}
 \boxed{
 1-u(t)
 =
 \left[
 \frac{Z_1\Tb^\beta(t)\Jn(t)}{\left[2\sigma\Jfour(t)\right]^2}
 \right]^{1/2}
 u^{(n+\beta)/8}(t).}
 \label{eq:u_equation}
\end{equation}
This equation is the only extra operation relative to the original surface-HR
formula.

\section{Bath-driven front formula}

Once $u(t)$ is known, the bath-driven front is
\begin{equation}
 \boxed{
 \xf^2(t)=K_\eps C\,
 u^{(n-\beta)/4}(t)\,\Tb^{-\beta}(t)\,\Jn(t).}
 \label{eq:bath_front}
\end{equation}
Equivalently,
\begin{equation}
 \boxed{
 \xf(t)=\left[K_\eps C\,
 u^{(n-\beta)/4}(t)\,\Tb^{-\beta}(t)\,\Jn(t)\right]^{1/2}.}
 \label{eq:bath_front_sqrt}
\end{equation}
If $u\to 1$, then $\Ts\to\Tb$ and \eqref{eq:bath_front} reduces to the naive
surface-HR formula with $\Ts=\Tb$.

\section{Temperature profile}

After computing $u(t)$ and $\xf(t)$, set
\begin{equation}
 \boxed{\Ts(t)=\Tb(t)u^{1/4}(t),\qquad H(t)=\Ts^n(t).}
\end{equation}
Define
\begin{equation}
 y=\frac{x}{\xf(t)}.
\end{equation}
The first-order HR profile is
\begin{equation}
 \boxed{
 T(x,t)=\Ts(t)
 \left[(1-y)\left(1+\frac{\eps}{2}(1-A)y\right)\right]^{1/[n(1-\eps)]}}
 \label{eq:profile}
\end{equation}
for $0\le y\le 1$, and $T=0$ for $y>1$.  The profile parameter is
\begin{equation}
 \boxed{
 A(t)=\frac{\xf^2(t)}{C H^{1-\eps}(t)}\frac{\dot H(t)}{H(t)}.}
 \label{eq:A}
\end{equation}
For a leading-order profile, use
\begin{equation}
 \boxed{
 T(x,t)\approx \Ts(t)\left(1-\frac{x}{\xf(t)}\right)^{1/n}.}
 \label{eq:leading_profile}
\end{equation}

\section{Numerical algorithm}

Given tabulated $\Tb(t)$:
\begin{enumerate}
 \item Interpolate $\Tb(t)$ on a convenient grid.  For noisy data, linear interpolation is usually safest.
 \item Compute $\Jfour(t_i)$ and $\Jn(t_i)$ by trapezoidal integration.
 \item For each time $t_i$, if $T_{b,i}=0$ or either fluence is zero, set $T_{s,i}=0$ and $x_{F,i}=0$.
 \item Otherwise solve \eqref{eq:u_equation} for $u_i\in(0,1)$.  Define
 \begin{equation}
  \Phi(u;t_i)=1-u-
  \left[\frac{Z_1 T_{b,i}^\beta J_{n,i}}{(2\sigma J_{4,i})^2}\right]^{1/2}
  u^{(n+\beta)/8}.
 \end{equation}
 Since $\Phi(0;t_i)>0$, $\Phi(1;t_i)<0$, and $\partial\Phi/\partial u<0$,
 there is a unique physical root.  Bisection or Brent on $(0,1)$ is robust.
 \item Compute $T_{s,i}=T_{b,i} u_i^{1/4}$.
 \item Compute $x_{F,i}$ from \eqref{eq:bath_front_sqrt}.
 \item For profiles, compute $H_i=T_{s,i}^n$, estimate $\dot H_i$ by finite differences, compute $A_i$ from \eqref{eq:A}, and use \eqref{eq:profile}.
\end{enumerate}

\section{Self-contained summary}

The bath-driven LTE HR model solves the diffusion equation
\begin{equation}
 \frac{\partial T^\beta}{\partial t}=C\frac{\partial^2T^n}{\partial x^2},
 \qquad
 n=4+\alpha,
 \qquad
 \eps=\frac{\beta}{n},
\end{equation}
with the material model
\begin{equation}
 e=fT^\beta\rho^{-\mu},
 \qquad
 \frac{1}{\kappa_R}=gT^\alpha\rho^{-\lambda},
\end{equation}
and the bath boundary condition
\begin{equation}
 F=2\sigma(\Tb^4-\Ts^4).
\end{equation}
Given $\Tb(t)$, compute
\begin{equation}
 \Jfour=\int_0^t\Tb^4(t')\dif t',
 \qquad
 \Jn=\int_0^t\Tb^n(t')\dif t'.
\end{equation}
Then solve
\begin{equation}
 1-u=\left[\frac{Z_1\Tb^\beta\Jn}{(2\sigma\Jfour)^2}\right]^{1/2}
 u^{(n+\beta)/8},
 \qquad
 Z_1=f^2\rho^{2(1-\mu)}(2+\eps)(1-\eps)C.
\end{equation}
The surface temperature is
\begin{equation}
 \Ts=\Tb u^{1/4},
\end{equation}
and the front is
\begin{equation}
 \boxed{
 \xf^2=\frac{2+\eps}{1-\eps}C\,
 u^{(n-\beta)/4}\Tb^{-\beta}\int_0^t\Tb^n(t')\dif t'.}
\end{equation}
The profile is
\begin{equation}
 \boxed{
 T(x,t)=\Ts(t)
 \left[(1-y)\left(1+\frac{\eps}{2}(1-A)y\right)\right]^{1/[n(1-\eps)]},
 \quad y=\frac{x}{\xf(t)},}
\end{equation}
with
\begin{equation}
 A=\frac{\xf^2}{C\Ts^{n(1-\eps)}}\frac{d}{dt}\ln \Ts^n.
\end{equation}

The HR spatial approximation has the same formal order as the original
Hammer--Rosen theory: the retained front and profile are first order in
$\eps$, with neglected HR spatial terms of order $O(\eps^2)$.  The additional
bath approximation is the fluence-ratio assumption
\begin{equation}
 \Ts(t')/\Tb(t')\approx \Ts(t)/\Tb(t)
 \quad (0<t'<t).
\end{equation}
Thus the model is most accurate when the surface-to-bath ratio varies smoothly
over the accumulated fluence history.  It is not the exact Cohen--Heizler
boundary evolution, but it is stable, compact, and avoids differentiating the
bath drive.

\end{document}
